\documentclass[11pt]{article}

\usepackage[margin=1in]{geometry}
\usepackage[T1]{fontenc}
\usepackage[utf8]{inputenc}
\usepackage{lmodern}
\usepackage{microtype}
\usepackage{booktabs}
\usepackage{longtable}
\usepackage{array}
\usepackage{xcolor}
\usepackage{hyperref}
\usepackage{enumitem}
\usepackage{listings}
\usepackage{amsmath,amssymb}
\usepackage{fancyhdr}
\usepackage{lastpage}

\hypersetup{
  colorlinks=true,
  linkcolor=blue!55!black,
  urlcolor=blue!55!black,
  citecolor=blue!55!black
}

\pagestyle{fancy}
\fancyhf{}
\lhead{Lean-AHR v2.0 Milestone}
\rhead{Benjamin A. Aenlle Falkenburg}
\cfoot{\thepage\ of \pageref{LastPage}}
\setlength{\headheight}{14pt}

\setlist[itemize]{topsep=2pt,itemsep=2pt,leftmargin=1.4em}
\setlist[enumerate]{topsep=2pt,itemsep=2pt,leftmargin=1.6em}
\renewcommand{\arraystretch}{1.15}

\lstdefinestyle{leanstyle}{
  basicstyle=\ttfamily\small,
  breaklines=true,
  frame=single,
  framerule=0.2pt,
  columns=fullflexible,
  keepspaces=true,
  backgroundcolor=\color{black!2}
}

\title{\textbf{Lean-AHR v2.0 Arithmetic Primitives Milestone Report}\\
\large Verified Proof Search, Arithmetic Bridge, and MVP Interaction Upgrade}
\author{Benjamin A. Aenlle Falkenburg}
\date{2026-05-31 Europe/Dublin}

\begin{document}
\maketitle

\begin{abstract}
This report documents the Lean-AHR v2.0 arithmetic-primitives milestone completed on 2026-05-31. The run
\texttt{lean\_ahr\_v2\_0\_arith\_primitives\_repeats27} evaluated 810 Lean-verified proof-search records across
top-1, top-2, and top-4 search budgets. It reached 100\% proof completion at every tested top-k level, including
the new v2.0 families for natural-number inequalities, divisibility, parity encoded as divisibility by 2, and
existential witness construction. The run produced a diagnostic atlas, milestone package, JSONL proof records,
checkmark logs, and an updated MVP command-line interface that can answer new arithmetic prompts and emit
Lean-checked proof files. This is a narrow but meaningful proof-search landmark: AHR is no longer only handling
propositional micro-proofs and addition rewrites; it now has a verified bridge into small arithmetic reasoning
patterns.
\end{abstract}

\tableofcontents
\newpage

\section{Executive Summary}

Lean-AHR v2.0 is the first milestone where the proof-search system consistently solved a mixed suite containing
both previous propositional/rewriting tasks and new arithmetic primitives. The important result is not merely that
the examples compiled in Lean. The important result is that the AHR-guided search loop selected tactic sequences,
validated local proof-state transitions with Lean, produced final theorem files, and completed the entire evaluation
without failures.

\begin{center}
\begin{tabular}{ll}
\toprule
Run name & \texttt{lean\_ahr\_v2\_0\_arith\_primitives\_repeats27} \\
Remote host at completion & \texttt{benjamin@35.239.179.142} \\
Project path & \texttt{/home/benjamin/03\_ahr\_lean\_foundation\_plan\_20260508} \\
Records & 810 \\
Top-k budgets & 1, 2, 4 \\
Completion & 100\% at top-1, top-2, and top-4 \\
Failure count & 0 \\
Clinical use flag & \texttt{forbidden} \\
GitHub archive commit & \texttt{f0ed3c1} \\
\bottomrule
\end{tabular}
\end{center}

The immediate interpretation is:

\begin{itemize}
  \item AHR's Lean proof-search layer is stable on this finite arithmetic/propositional benchmark.
  \item The top-1 policy is strong on v2.0: it solved all 270 top-1 records.
  \item The wider search budgets did not expose hidden brittleness: top-2 and top-4 also reached 100\%.
  \item The MVP interface is now able to turn a small set of controlled English prompts into Lean theorem templates and verified proofs.
  \item The next scientific step is not a bigger demo claim; it is less-templated generalization in v2.1.
\end{itemize}

\section{What Lean-AHR Is Doing}

Lean-AHR is being built as a verified reasoning layer around Lean. The working model is:

\begin{enumerate}
  \item A theorem statement defines a formal target.
  \item The current proof state is treated as a finite decision context.
  \item Candidate tactics are generated from the local state, theorem family, hypotheses, target shape, and known proof patterns.
  \item AHR scores candidate actions using a learned branch-cost/ranking model plus hand-coded safety priors.
  \item Lean checks each local step. Invalid branches are rejected by the verifier rather than trusted.
  \item If a complete proof is found, a final Lean file is written and compiled without \texttt{sorry}.
\end{enumerate}

This architecture deliberately avoids treating the neural model as an oracle. The neural component is used to rank
finite choices. Lean is the final source of truth. In other words, AHR proposes; Lean disposes.

\section{Mathematical Scope of v2.0}

Before v2.0, the strongest Lean-AHR results were primarily in propositional logic, local lemma construction, cases,
contradiction, and small natural-number addition rewrites such as associativity and commutativity. v2.0 extends the
mathematical surface area to the following natural-number primitives.

\subsection{Inequality}

The v2.0 inequality family included reflexivity, zero lower bounds, successor bounds, additive bounds, and
transitivity:

\begin{align*}
n &\le n,\\
0 &\le n,\\
n &< n.\mathrm{succ},\\
n &\le n + m,\\
n &\le m + n,\\
n \le m \wedge m \le k &\Rightarrow n \le k.
\end{align*}

The key Lean tactics included direct use of \texttt{Nat.le\_refl}, \texttt{Nat.le\_add\_right},
\texttt{Nat.le\_add\_left}, and \texttt{Nat.le\_trans}. This is modest arithmetic, but it matters because it
introduces order reasoning and hypothesis chaining beyond equality rewriting.

\subsection{Divisibility}

The divisibility family introduced goals of the form:

\begin{align*}
n &\mid n,\\
n \mid m \wedge n \mid k &\Rightarrow n \mid m+k,\\
n \mid m \wedge m \mid k &\Rightarrow n \mid k.
\end{align*}

The proof search used Lean primitives such as \texttt{Nat.dvd\_refl}, \texttt{Nat.dvd\_add}, and
\texttt{Nat.dvd\_trans}. This extends AHR into relational arithmetic where proof success depends on selecting the
right hypothesis pair, not only rewriting a target.

\subsection{Parity as Divisibility}

Lean core did not expose a convenient \texttt{Even}/\texttt{Odd} layer in the current no-mathlib setup, so v2.0
encoded parity through divisibility by 2:

\[
2 \mid n \Rightarrow 2 \mid n+n.
\]

This was solved using \texttt{Nat.dvd\_add h h}. It is a useful engineering compromise: it keeps the benchmark inside
core Lean while still representing a common arithmetic concept.

\subsection{Existential Witnesses}

The existential family introduced witness construction:

\[
\exists k : \mathbb{N},\ k = n.
\]

The proof term is short but structurally important:

\begin{lstlisting}[style=leanstyle]
exact Exists.intro n rfl
\end{lstlisting}

This matters because future theorem proving cannot only discharge universally quantified equations. It must also
construct objects. v2.0 adds the first reliable witness-construction path.

\section{Engineering Changes}

\subsection{New v2.0 Runner}

The main experiment script was:

\begin{lstlisting}[style=leanstyle]
tools/run_lean_ahr_v2_0_arith_primitives.py
\end{lstlisting}

The detached launcher was:

\begin{lstlisting}[style=leanstyle]
tools/run_lean_ahr_v2_0_arith_primitives_detached.sh
\end{lstlisting}

The run was explicitly resumable. This mattered operationally because the VM/IP changed during the job. The run had
already written 607 row-level records out of 810. After reconnecting to \texttt{35.239.179.142}, the runner resumed
from the existing JSONL file rather than restarting the whole experiment. It then completed the remaining top-4
records and wrote the final report.

\subsection{Detached Execution and Evidence Trail}

The execution pattern produced:

\begin{itemize}
  \item a PID file,
  \item detached stdout/stderr log,
  \item GPU telemetry log,
  \item 15-minute checkmark snapshots,
  \item row-level JSONL proof records,
  \item progress JSON,
  \item final report JSON,
  \item diagnostic atlas,
  \item milestone package with SHA-256 manifest.
\end{itemize}

This is important for research discipline. A successful proof-search run is not just a console message; it has
replayable artifacts, row-level records, and a packaged audit trail.

\subsection{MVP CLI Upgrade}

The interactive MVP CLI was upgraded:

\begin{lstlisting}[style=leanstyle]
tools/lean_ahr_mvp_cli.py
\end{lstlisting}

It now supports prompt templates and curated examples for:

\begin{itemize}
  \item addition associativity and commutativity,
  \item addition reordering,
  \item implication chains,
  \item conjunction projection,
  \item contradiction,
  \item disjunction/case-style bridges,
  \item natural-number inequalities,
  \item natural-number divisibility,
  \item parity as divisibility by 2,
  \item existential witness construction.
\end{itemize}

The smoke suite:

\begin{lstlisting}[style=leanstyle]
tools/run_lean_ahr_mvp_smoke_suite.sh
\end{lstlisting}

was expanded to include v2.0 arithmetic prompts and solved all six tested examples.

\section{Experiment Results}

\subsection{Top-k Completion}

\begin{center}
\begin{tabular}{rrrrrr}
\toprule
Top-k & Count & Solved & Completion & Mean Lean calls & Mean elapsed s \\
\midrule
1 & 270 & 270 & 1.000 & 16.20 & 6.46 \\
2 & 270 & 270 & 1.000 & 35.93 & 13.68 \\
4 & 270 & 270 & 1.000 & 64.53 & 25.32 \\
\bottomrule
\end{tabular}
\end{center}

The result is unusually clean: all tested search budgets passed. The top-1 score is especially useful because it
means the ranker and tactic priors are now choosing the correct first branch across the v2.0 suite, not merely
recovering through wider search.

\subsection{Acceptance Gates}

\begin{center}
\begin{tabular}{ll}
\toprule
Gate & Result \\
\midrule
top-1 completion $\ge 0.90$ & pass \\
top-2 completion $\ge 0.90$ & pass \\
top-4 completion $\ge 0.95$ & pass \\
zero final Lean failures & pass \\
clinical use marked forbidden & pass \\
\bottomrule
\end{tabular}
\end{center}

\subsection{Family-Level Coverage}

The v2.0 run intentionally retained prior proof families while adding arithmetic primitives. This tests regression as
well as new capability.

\begin{longtable}{p{0.42\linewidth}rrrr}
\toprule
Family & Count & Solved & Completion & Mean Lean calls \\
\midrule
\endhead
\texttt{have\_construct} & 27 & 27 & 1.000 & 86.00 \\
\texttt{have\_contradiction} & 27 & 27 & 1.000 & 22.67 \\
\texttt{have\_imp\_chain} & 27 & 27 & 1.000 & 48.67 \\
\texttt{have\_or\_elim} & 27 & 27 & 1.000 & 46.67 \\
\texttt{have\_projection} & 27 & 27 & 1.000 & 36.00 \\
\texttt{v16\_deep\_chain\_conj} & 27 & 27 & 1.000 & 139.00 \\
\texttt{v16\_nested\_or\_bridge} & 27 & 27 & 1.000 & 56.33 \\
\texttt{v16\_or\_to\_local\_conj} & 27 & 27 & 1.000 & 72.67 \\
\texttt{v16\_project\_contradiction} & 27 & 27 & 1.000 & 57.67 \\
\texttt{v17\_contra\_with\_distractor} & 27 & 27 & 1.000 & 83.67 \\
\texttt{v17\_deep\_reordered\_bridge} & 27 & 27 & 1.000 & 84.33 \\
\texttt{v17\_distractor\_projection} & 27 & 27 & 1.000 & 47.67 \\
\texttt{v17\_or\_with\_distractor} & 27 & 27 & 1.000 & 46.67 \\
\texttt{v17\_reordered\_chain} & 27 & 27 & 1.000 & 52.67 \\
\texttt{v18\_nat\_assoc} & 27 & 27 & 1.000 & 9.33 \\
\texttt{v18\_nat\_comm} & 27 & 27 & 1.000 & 11.67 \\
\texttt{v18\_nat\_reorder} & 27 & 27 & 1.000 & 58.00 \\
\texttt{v18\_nat\_rewrite} & 54 & 54 & 1.000 & 10.17 \\
\texttt{v20\_nat\_divisibility} & 81 & 81 & 1.000 & 23.89 \\
\texttt{v20\_nat\_exists} & 27 & 27 & 1.000 & 8.33 \\
\texttt{v20\_nat\_inequality} & 135 & 135 & 1.000 & 10.40 \\
\texttt{v20\_nat\_inequality\_trans} & 27 & 27 & 1.000 & 26.67 \\
\texttt{v20\_nat\_parity\_div2} & 27 & 27 & 1.000 & 28.67 \\
\bottomrule
\end{longtable}

\section{MVP Smoke Results}

After the proof-search run passed, the MVP CLI was updated and tested. The smoke suite solved six out of six examples.

\begin{center}
\small
\begin{tabular}{p{0.32\linewidth}p{0.40\linewidth}rr}
\toprule
Prompt or task & Proof sketch & Lean calls & Solved \\
\midrule
prove addition is associative & \texttt{rw [Nat.add\_assoc]} & 4 & yes \\
prove a natural number is less equal itself & \texttt{exact Nat.le\_refl n} & 4 & yes \\
prove divisibility addition & \texttt{intro h0; intro h1; exact Nat.dvd\_add h0 h1} & 11 & yes \\
prove parity of a double & \texttt{intro h0; exact Nat.dvd\_add h0 h0} & 8 & yes \\
prove existential witness identity & \texttt{exact Exists.intro n rfl} & 4 & yes \\
custom projection theorem & \texttt{intro h0; intro h1; exact h0} & 31 & yes \\
\bottomrule
\end{tabular}
\end{center}

The MVP is still a controlled prototype, not an open-ended theorem prover. However, it is now interactive enough to
serve as a real demonstration: a user can give a supported prompt or theorem statement and receive a Lean-checked
proof file.

\section{Artifact Inventory}

The milestone package is archived locally and pushed to GitHub.

\begin{lstlisting}[style=leanstyle]
Local archive:
/Users/benjaminfalkenburg/Documents/primumai/github_exports/
  ahr_verified_reasoning_20260509/
    remote_tools/remote_lean_35_239_179_142_20260531/
      v2_0_arith_primitives/
        lean_ahr_v2_0_arith_primitives_repeats27_MILESTONE_PACKAGE_20260531T221022Z/

GitHub:
https://github.com/BoredDuffa/ahr-verified-reasoning-20260509

Commit:
f0ed3c1 Archive Lean AHR v2.0 arithmetic milestone
\end{lstlisting}

The milestone package contains:

\begin{itemize}
  \item final report JSON,
  \item progress report JSON,
  \item 810-record JSONL proof trace dataset,
  \item detached logs,
  \item GPU telemetry logs,
  \item pre-resume logs,
  \item checkmark snapshots,
  \item diagnostic atlas JSON and Markdown,
  \item updated MVP CLI,
  \item MVP smoke suite and smoke results,
  \item SHA-256 manifest.
\end{itemize}

\section{Interpretation}

This milestone means AHR has crossed an important narrow threshold. It can now solve and verify a controlled set of
Lean theorem-proving tasks that require:

\begin{itemize}
  \item proposition introduction,
  \item projection from conjunctions,
  \item contradiction handling,
  \item disjunction/case-style bridges,
  \item local lemma construction,
  \item addition rewriting,
  \item inequality facts,
  \item divisibility facts,
  \item witness construction for existential goals.
\end{itemize}

The system is not yet proving advanced number theory. It is not yet proving irrationality of $\sqrt{2}$. But the
engineering direction is now clearer: build outward from verified local theorem families, keep Lean as the final
verifier, and only claim capability when a held-out proof-search run and an interactive MVP smoke both pass.

\section{Current Limitations}

\begin{itemize}
  \item The theorem templates are still narrow. v2.0 success does not prove broad theorem-proving generality.
  \item The prompt translator is a rule-based template router, not a trained language model.
  \item The current benchmark mostly uses Lean core. Mathlib-backed theorem proving is a separate next step.
  \item The system is CPU/Lean-verifier bound during proof search; GPU utilization is naturally low for these runs.
  \item Search completeness is bounded by max depth, max nodes, tactic candidates, and timeout.
  \item The benchmark contains validation/test-style splits, but v2.1 needs more aggressive de-templating to test real generalization.
\end{itemize}

\section{Next Research Step: v2.1}

The correct next step is Lean-AHR v2.1: less-templated arithmetic generalization.

\subsection{Recommended v2.1 Families}

\begin{enumerate}
  \item Reordered inequality transitivity:
  \[
  m \le k \rightarrow n \le m \rightarrow n \le k.
  \]
  \item Divisibility chains with distractor hypotheses:
  \[
  n \mid m \rightarrow x \mid y \rightarrow m \mid k \rightarrow n \mid k.
  \]
  \item Existential witnesses with nontrivial witness terms:
  \[
  \exists k : \mathbb{N},\ k = n + 0.
  \]
  \item Mixed propositional and arithmetic gates:
  \[
  P \rightarrow (P \rightarrow n \le m) \rightarrow n \le m.
  \]
  \item Parity/divisibility goals with irrelevant assumptions.
\end{enumerate}

\subsection{v2.1 Acceptance}

The proposed gates should be:

\begin{itemize}
  \item top-1 completion at least 90\%,
  \item top-2 completion at least 95\%,
  \item top-4 completion at least 98\%,
  \item zero Lean expectation mismatches,
  \item all final proofs compile without \texttt{sorry},
  \item a failure atlas for every unsolved theorem.
\end{itemize}

\section{Path Toward Mathlib and Harder Proofs}

The route from v2.0 to proofs such as irrationality of $\sqrt{2}$ should be staged.

\begin{enumerate}
  \item Finish v2.1 less-templated arithmetic.
  \item Add a Mathlib environment smoke test and import discipline.
  \item Build a small Mathlib theorem selection set: divisibility lemmas, parity lemmas, prime facts, square facts.
  \item Create a retrieval layer for named lemmas.
  \item Evaluate local proof search where AHR chooses among lemma applications.
  \item Only then attempt multi-lemma proofs such as irrationality of $\sqrt{2}$.
\end{enumerate}

The mathematical challenge in $\sqrt{2}$ is not simply arithmetic. It requires a chain of definitions and lemmas about
irrationality, divisibility, parity, squares, primes, and contradiction. Lean can verify such proofs, but AHR needs a
controlled way to discover the required lemmas and route through them.

\section{Conclusion}

Lean-AHR v2.0 is a clean milestone. The run completed 810 Lean-verified proof-search records with 100\% completion
at top-1, top-2, and top-4. It introduced working arithmetic primitives, preserved earlier propositional capability,
upgraded the MVP interaction path, and produced a packaged, checksummed artifact set pushed to GitHub.

The next milestone should deliberately stress generalization. v2.1 should keep the same verifier discipline but make
the arithmetic proofs less templated, more reordered, and more distractor-heavy. If that passes, the project will be
ready for its first serious Mathlib bridge.

\appendix
\section{Handoff Commands}

\subsection{Check Remote Status}

\begin{lstlisting}[style=leanstyle]
ssh benjamin@35.239.179.142
cd /home/benjamin/03_ahr_lean_foundation_plan_20260508
ls -lh reports/lean_ahr_v2_0_arith_primitives_repeats27_report.json
tail -n 40 reports/lean_ahr_v2_0_arith_primitives_repeats27.detached.log
\end{lstlisting}

\subsection{Run MVP Examples}

\begin{lstlisting}[style=leanstyle]
cd /home/benjamin/03_ahr_lean_foundation_plan_20260508
/opt/conda/bin/python tools/lean_ahr_mvp_cli.py --list-examples
/opt/conda/bin/python tools/lean_ahr_mvp_cli.py \
  --prompt "prove divisibility addition" \
  --name demo_dvd_add \
  --run-name lean_ahr_mvp_demo_dvd_add \
  --max-k 1 \
  --cpu
\end{lstlisting}

\subsection{Archive Location}

\begin{lstlisting}[style=leanstyle]
/Users/benjaminfalkenburg/Documents/primumai/github_exports/
  ahr_verified_reasoning_20260509/
    remote_tools/remote_lean_35_239_179_142_20260531/
\end{lstlisting}

\subsection{GitHub Push Already Completed}

\begin{lstlisting}[style=leanstyle]
cd /Users/benjaminfalkenburg/Documents/primumai/github_exports/ahr_verified_reasoning_20260509
git log --oneline -1
# f0ed3c1 Archive Lean AHR v2.0 arithmetic milestone
\end{lstlisting}

\end{document}
